%\setcounter{framenumber}{30}


\renewcommand{\mydate}{2025年03月18号\quad 第七次课}

\makeatletter
\ifodd\value{page}
\else
\begin{frame}[plain]
\end{frame}
\fi
\makeatother


\section{行列式}


\title{第四章 \quad 行列式}
\author{}
\date{}


\frame{\titlepage}


\subsection{行列式定义}



\subsubsection{行列式的引入-解线性方程组}
\begin{frame}\frametitle{行列式的引入}
考察线性方程组
\begin{equation} \tag{$*$}
\left\{\begin{array}{ccccccccc}
a_{11}x_1& +& a_{12}x_2 & + & \cdots & + & a_{1n}x_n &= & b_1\\
a_{21}x_1& +& a_{22}x_2 & + & \cdots & + & a_{2n}x_n &= & b_2\\
\vdots &&  \vdots &&  \ddots &&  \vdots &&  \vdots \\
a_{m1}x_1& +& a_{m2}x_2 & + & \cdots & + & a_{mn}x_n &= & b_m\\
\end{array}\right.
\end{equation}
\pp
\bigskip

当$n=2$时,若$a_{11}a_{22}-a_{12}a_{21}\neq0$, 则通过消元法可证明方程组有唯一解
\begin{equation*}
\left\{\begin{array}{ccccccccc}
x_1 &=& \frac{b_1a_{22}-b_2a_{12}}{a_{11}a_{22}-a_{12}a_{21}} \\
x_2 &=& \frac{b_2a_{11}-b_1a_{21}}{a_{11}a_{22}-a_{12}a_{21}} \\
\end{array}\right.
\end{equation*}
\pp

对任意四个数 $a,b,c,d$, 引入记号
\[\det\begin{pmatrix}a&b\\c&d\\\end{pmatrix} := \begin{vmatrix}a&b\\c&d\\\end{vmatrix} := ad - bc,\]
则方程的公式解可简洁地表示为
\[x_1 = \frac{\begin{vmatrix}b_1&a_{12}\\b_2&a_{22}\\\end{vmatrix}}{\begin{vmatrix}a_{11}&a_{12}\\a_{21}&a_{22}\\\end{vmatrix}}, \quad x_2 = \frac{\begin{vmatrix}a_{11}&b_{1}\\a_{21}&b_{2}\\\end{vmatrix}}{\begin{vmatrix}a_{11}&a_{12}\\a_{21}&a_{22}\\\end{vmatrix}}.\]
\end{frame}





\begin{frame}\frametitle{}
当$n=3$时, 方程组为
\begin{equation*}
\left\{\begin{array}{ccccccc}
a_{11}x_1& +& a_{12}x_2 & +  & a_{13}x_3 &= & b_1\\
a_{21}x_1& +& a_{22}x_2 & +  & a_{23}x_3 &= & b_2\\
a_{31}x_1& +& a_{32}x_2 & +  & a_{33}x_3 &= & b_3\\
\end{array}\right.
\end{equation*}
\pp
\bigskip

若方程组有唯一解,通过消元法(见课本p64-65)可得

\begin{equation*}\left\{
\begin{aligned}
x_1 &= \frac{
b_1\begin{vmatrix}a_{22}&a_{23}\\a_{32}&a_{33}\\\end{vmatrix}
-b_2\begin{vmatrix}a_{12}&a_{13}\\a_{32}&a_{33}\\\end{vmatrix}
+b_3\begin{vmatrix}a_{12}&a_{13}\\a_{22}&a_{23}\\\end{vmatrix}
}{
a_{11}\begin{vmatrix}a_{22}&a_{23}\\a_{32}&a_{33}\\\end{vmatrix}
-a_{21}\begin{vmatrix}a_{12}&a_{13}\\a_{32}&a_{33}\\\end{vmatrix}
+a_{31}\begin{vmatrix}a_{12}&a_{13}\\a_{22}&a_{23}\\\end{vmatrix}} \\[1mm]
x_2 &= \frac{
-b_1\begin{vmatrix}a_{21}&a_{23}\\a_{31}&a_{33}\\\end{vmatrix}
+b_2\begin{vmatrix}a_{11}&a_{13}\\a_{31}&a_{33}\\\end{vmatrix}
-b_3\begin{vmatrix}a_{11}&a_{13}\\a_{21}&a_{23}\\\end{vmatrix}
}{
a_{11}\begin{vmatrix}a_{22}&a_{23}\\a_{32}&a_{33}\\\end{vmatrix}
-a_{21}\begin{vmatrix}a_{12}&a_{13}\\a_{32}&a_{33}\\\end{vmatrix}
+a_{31}\begin{vmatrix}a_{12}&a_{13}\\a_{22}&a_{23}\\\end{vmatrix}} \\[1mm]
x_3 &= \frac{
b_1\begin{vmatrix}a_{21}&a_{22}\\a_{31}&a_{32}\\\end{vmatrix}
-b_2\begin{vmatrix}a_{11}&a_{12}\\a_{31}&a_{32}\\\end{vmatrix}
+b_3\begin{vmatrix}a_{11}&a_{12}\\a_{21}&a_{22}\\\end{vmatrix}
}{
a_{11}\begin{vmatrix}a_{22}&a_{23}\\a_{32}&a_{33}\\\end{vmatrix}
-a_{21}\begin{vmatrix}a_{12}&a_{13}\\a_{32}&a_{33}\\\end{vmatrix}
+a_{31}\begin{vmatrix}a_{12}&a_{13}\\a_{22}&a_{23}\\\end{vmatrix}}
\end{aligned}\right.
\end{equation*}
\end{frame}




\begin{frame}\frametitle{}
 引入记号
\begin{equation*}
\begin{aligned}
\det&\begin{pmatrix}a_{11}&a_{12}&a_{13}\\a_{21}&a_{22}&a_{23}\\a_{31}&a_{32}&a_{33}\\\end{pmatrix}
   := \begin{vmatrix}a_{11}&a_{12}&a_{13}\\a_{21}&a_{22}&a_{23}\\a_{31}&a_{32}&a_{33}\\\end{vmatrix}\\
 & := a_{11}\begin{vmatrix}a_{22}&a_{23}\\a_{32}&a_{33}\\\end{vmatrix}
     -a_{21}\begin{vmatrix}a_{12}&a_{13}\\a_{32}&a_{33}\\\end{vmatrix}
     +a_{31}\begin{vmatrix}a_{12}&a_{13}\\a_{22}&a_{23}\\\end{vmatrix}.\\
\end{aligned}
\end{equation*}
\pp
\bigskip

则方程的公式解可简化为
\begin{equation*}
x_1 = \frac{
\begin{vmatrix}b_{1}&a_{12}&a_{13}\\b_{2}&a_{22}&a_{23}\\b_{3}&a_{32}&a_{33}\\\end{vmatrix}
}{
\begin{vmatrix}a_{11}&a_{12}&a_{13}\\a_{21}&a_{22}&a_{23}\\a_{31}&a_{32}&a_{33}\\\end{vmatrix}}
\quad x_2 = \frac{
\begin{vmatrix}a_{11}&b_{1}&a_{13}\\a_{21}&b_{2}&a_{23}\\a_{31}&b_{3}&a_{33}\\\end{vmatrix}
}{
\begin{vmatrix}a_{11}&a_{12}&a_{13}\\a_{21}&a_{22}&a_{23}\\a_{31}&a_{32}&a_{33}\\\end{vmatrix}}
\quad x_3 = \frac{
\begin{vmatrix}a_{11}&a_{12}&b_{1}\\a_{21}&a_{22}&b_{2}\\a_{31}&a_{32}&b_{3}\\\end{vmatrix}
}{
\begin{vmatrix}a_{11}&a_{12}&a_{13}\\a_{21}&a_{22}&a_{23}\\a_{31}&a_{32}&a_{33}\\\end{vmatrix}}
\end{equation*}
\pp
\bigskip

利用归纳法, 我们可以引入$n$-阶行列式, 将其定义为$n$个$n-1$行列式的某个组合. 从而将一般的$n$-元线性方程组的公式解表示为两个 $n$-阶行列式的比值.
\end{frame}

\subsubsection{行列式的引入-体积}
\begin{frame}\frametitle{行列式的引入}
\begin{equation*}
\begin{tikzpicture}[scale=0.5]
\tikzmath{\r=8;};
\coordinate (O) at (0,0);
\coordinate (A) at (1,2);
\coordinate (B) at (1.5,-0.5);
\coordinate (C) at (2.5,1.5);
\draw[->,thick,blue](O)--(A);
\node [above] at (A) {$(c,d)$};
\draw[->,thick,blue](O)--(B);
\node [right] at (B) {$(a,b)$};
\draw[dotted](A)--(C);
\draw[dotted](B)--(C);
\coordinate (o) at (\r+0,0);
\coordinate (a) at (\r+1,-1);
\coordinate (b) at (\r+1,1);
\coordinate (c) at (\r+1,2);
\coordinate (ab) at (\r+2,0);
\coordinate (ac) at (\r+2,1);
\coordinate (bc) at (\r+2,3);
\coordinate (abc) at (\r+3,2);
\draw[->,thick,blue](o)--(a);\node [below] at (a) {\tiny $(a_{11},a_{12},a_{13})$};
\draw[->,dotted,thick,blue](o)--(b);\node [left] at (b) {\tiny $(a_{21},a_{22},a_{23})$};
\draw[->,thick,blue](o)--(c);\node [above left] at (c) {\tiny $(a_{31},a_{32},a_{33})$};
\draw[thick,blue](a)--(ab);
\draw[thick,blue](a)--(ac);
\draw[dotted,thick,blue](b)--(ab);
\draw[dotted,thick,blue](b)--(bc);
\draw[thick,blue](c)--(ac);
\draw[thick,blue](c)--(bc);
\draw[thick,blue](ab)--(abc);
\draw[thick,blue](ac)--(abc);
\draw[thick,blue](bc)--(abc);
\end{tikzpicture}
\end{equation*}
\pp
\begin{itemize}
\item 给定两个二维数组向量$(a,b)$和$(c,d)$.如图我们有平行四边形. \pp
这时平行四边形的(有向)面积为 $S=ad-bc=:\begin{vmatrix} a & b\\ c & d \\\end{vmatrix}$.
\pp
\item 给定三个三维数组向量$(a_{11},a_{12},a_{13})$,$(a_{21},a_{22},a_{23})$和$(a_{31},a_{32},a_{33})$.如图我们可构造一个平行六面体. \pp 其(有向)体积正好为
\[\Delta= a_{11}\begin{vmatrix} a_{22} & a_{23} \\ a_{32} & a_{33} \end{vmatrix}
- a_{12}\begin{vmatrix} a_{21} & a_{23} \\ a_{31} & a_{33} \end{vmatrix}
+ a_{13}\begin{vmatrix} a_{21} & a_{22} \\ a_{31} & a_{32} \end{vmatrix}.
\]
\item 若给定$n$个$n$维数组向量呢？
\end{itemize}

\end{frame}

% ========== 2023年03月28号 第07次课 行列式 ==========
%\frame{\titlepage}





\subsubsection{行列式的定义}
\begin{frame}\frametitle{行列式的定义(高维平行多面体的有向体积)}

\begin{definition}[行列式]
方阵$A=(a_{ij})_{n\times n}$的\blue{行列式}记为
\[\det(A) \quad \text{或} \quad
\begin{vmatrix}
a_{11}&a_{12}&\cdots&a_{1n}\\
a_{21}&a_{22}&\cdots&a_{2n}\\
\cdots&\cdots&\cdots&\cdots\\
a_{n1}&a_{n2}&\cdots&a_{nn}\\
\end{vmatrix}_.\]
\pp
当$n=1$时,
\[\det(A):=a_{11}.\]
\pp
当$n\geq2$时, $\det(A)$递归地定义为
\[\det(A):=\sum_{k=1}^n(-1)^{k+1}a_{k1}\cdot
\begin{vmatrix}
a_{12}&a_{13}&\cdots&a_{1n}\\
\cdots&\cdots&\cdots&\cdots\\
a_{k-1,2}&a_{k-1,3}&\cdots&a_{k-1,n}\\
a_{k+1,2}&a_{k+1,3}&\cdots&a_{k+1,n}\\
\cdots&\cdots&\cdots&\cdots\\
a_{n2}&a_{n,3}&\cdots&a_{nn}\\
\end{vmatrix} \]
\end{definition}
\end{frame}





\subsubsection{从定义出发计算行列式的可行性分析}
\begin{frame}\frametitle{行列式}
\begin{example}
$\begin{vmatrix}a&b\\c&d\end{vmatrix}=ad-bc$,
\pp
$\begin{vmatrix}a_1&a_2&a_3\\b_1&b_2&b_3\\c_1&c_2&c_3\end{vmatrix}=a_1b_2c_3+a_2b_3c_1+a_3b_1c_2-a_1b_3c_2-a_2b_1c_3-a_3b_2c_1$.
\end{example}
\pp
\begin{question} 用定义计算$n$阶行列式需要计算多少个乘法?
\end{question}
\pp
\[10 \text{ 阶行列式} \ \rightsquigarrow \ 10!\times 9 \text{ 个乘法} \ \rightsquigarrow \ \frac{10!\times 9}{24\times60\times60}=378\text{天}.\]
\end{frame}





\subsubsection{子矩阵}
\begin{frame}\frametitle{子矩阵}
\pp
\begin{definition}
由矩阵$A=(a_{ij})_{m\times n}$的第$i_1,i_2,\cdots,i_r$行和第$j_1,j_2,\cdots,j_s$列上的元素依次排列组成的$r\times s$矩阵称为$A$的\blue{子矩阵},记作
\[A\left({i_1i_2\cdots i_r\atop j_1j_2\cdots j_s}\right) = \begin{pmatrix}
a_{i_1j_1}&a_{i_1j_2}&\cdots&a_{i_1j_s}\\
a_{i_2j_1}&a_{i_2j_2}&\cdots&a_{i_2j_s}\\
\cdots&\cdots&\cdots&\cdots\\
a_{i_rj_1}&a_{i_rj_2}&\cdots&a_{i_rj_s}\\
\end{pmatrix}\]
\end{definition}
\pp
\begin{example}
若 $A=\left(\begin{array}{cccc}
11 & 12 & 13 & 14\\
21 & 22 & 23 & 24\\
31 & 32 & 33 & 34\\
41 & 42 & 43 & 44\\
\end{array}\right)$,
\pp
则 $A\left({{1 3}\atop {1 3}}\right)  = \pp \left(\begin{array}{cccc}
11 & 13 \\
31 & 33 \\
\end{array}\right)$;
\pp
$A\left({{412}\atop {13}}\right) = \pp \left(\begin{array}{cccc}
41 & 43\\
11 & 13\\
21 & 23\\
\end{array}\right)$;
\end{example}
\end{frame}





\subsubsection{代数余子式}
\begin{frame}\frametitle{代数余子式}
\begin{definition}
给定一个$n$阶矩阵$A=(a_{ij})_{n\times n}$.
\pp
\begin{enumerate}
\item 称 $M_{ij}:=\det\left(A\left({{1,\cdots,i-1,i+1,\cdots,n}\atop{1,\cdots,j-1,j+1,\cdots,n}}\right)\right)$为元素$a_{ij}$的\blue{余子式}.
\pp
\item 称$A_{ij}:=(-1)^{i+j}M_{ij}$为$a_{ij}$的\blue{代数余子式}.
\pp
\item 更一般地,称$A$的某个$k$阶子矩阵的行列式为$A$的\blue{$k$阶子式},删去子式所在的行列所得的矩阵的行列式称为该子式的\blue{余子式}.
\pp
\item 特别地,我们称$\det\left(A\left({{1,2,\cdots,k}\atop{1,2,\cdots,k}}\right)\right)$为$A$的\blue{$k$阶主子式}.
\end{enumerate}
\end{definition}
\pp
根据这个定义,我们可以将行列式的定义式简写为
\[\boxed{\det(A) := \sum_{i=1}^na_{i1}A_{i1} = \sum_{i=1}^{n}(-1)^{i+1}a_{i1}M_{i1}.}\]
\end{frame}






\begin{frame}\frametitle{行列式例子}
\begin{example}
\[\begin{vmatrix}
a_{11}&a_{12}&\cdots&a_{1n}\\
  &a_{22}&\ddots&\vdots\\
  &      &\ddots&a_{n2}\\
  &      &      &a_{nn}\\
\end{vmatrix}=a_{11}a_{22}\cdots a_{nn}.\]
\end{example}
\begin{example}
\[\begin{vmatrix}
  &      &      &a_{1n}\\
  &      &\iddots&a_{2n}\\
  & a_{n-1,2}&\iddots&\vdots\\
a_{n1} & a_{n2} & \cdots & a_{nn}\\
\end{vmatrix}=(-1)^{\frac{n(n-1)}{2}}a_{n1}a_{n-1,2}\cdots a_{1n}.\]
\end{example}
\end{frame}


\subsection{行列式性质}


\subsubsection{行列式列展开}
\begin{frame}\frametitle{行列式列展开}
\begin{theorem}[行列式列展开]
设$A=(a_{ij})_{n\times n}$为$n$阶方阵.则对任意$1\leq j \leq n$
\[\boxed{\det(A)=\sum_{i=1}^na_{ij}A_{ij} = \sum_{i=1}^{n}(-1)^{i+j}a_{ij}M_{ij}.}\]
\end{theorem}
\yang{思路: 根据对行列式阶数的归纳, 将行列式
\[\det(A):=\sum_{i'=1}^{n}(-1)^{i'+1}a_{i'1}M_{i'1}\]
中的$M_{i'1}$ 按第 $j-1$列(即,原行列式的第$j$行)再次展开. 同时右式
\[\sum_{i=1}^{n}(-1)^{i+j}a_{ij}M_{ij}\]
中的 $M_{ij}$ 按第一列展开. 然后比较两式是否相等即可.}
\end{frame}





\begin{frame}\frametitle{证明细节}
{\tiny 记 $M\left[{{i_1i_2}\atop {j_1j_2}}\right]$ 为矩阵$A$去掉第$i_1,i_2$行和$j_1,j_2$列后的行列式。即
\[ M\left[{{i_1i_2}\atop {j_1j_2}}\right]:=\det\left(A\left({1,\cdots,i_1-1,i_1+1,\cdots,i_2-1,i_2+1,\cdots,n} \atop {1,\cdots,j_1-1,j_1+1,\cdots,j_2-1,j_2+1,\cdots,n}\right)\right).\]
\pp
对行列式的阶数$n$进行归纳。当$n=1$,结论显然成立。假设结论对n-1阶行列式成立。
\pp
则}
\begin{equation*} \tiny
\begin{split}
\det(A)
& = \sum_{i' =1}^{n} (-1)^{i' +1}a_{i' 1}M_{i' 1} \\ \pp
& = \sum_{i' =1}^{n} (-1)^{i' +1}a_{i' 1} \left(\sum_{i =1}^{i' -1} (-1)^{(j-1)+i } a_{i j}  \cdot M\left[{i' ,i }\atop{1,j}\right] + \sum_{i =i' +1}^{n}(-1)^{(i -1)+(j-1)} a_{i j} \cdot M\left[{i' ,i }\atop{1,j}\right] \right)\\ \pp
& = \sum_{i =1}^{n} (-1)^{i +j}a_{i j}\left(\sum_{i' =1}^{i -1} (-1)^{i' +1} a_{i' 1} M\left[{i' ,i }\atop{1,j}\right] + \sum_{i' =i +1}^{n}(-1)^{(i' -1)+1} a_{i' 1}  \cdot  M\left[{i' ,i }\atop{1,j}\right] \right)\\ \pp
& = \sum_{i =1}^{n} (-1)^{i +j}a_{i j} M_{i j} \\  \pp
& = \sum_{i=1}^{n} a_{ij}A_{ij} \\ \pp
\end{split}
\end{equation*}
\pp
{\tiny 其中第二个等式为对$n-1$阶行列式$M_{i'1}$的第$j-1$列进行展开，第三个等式为交换求和号，第四个等式为对$n-1$阶行列式$M_{ij}$的第$1$列进行展开。
}
\end{frame}



\begin{frame}\frametitle{行列式例子}
\begin{example}
\[\begin{vmatrix}
a_{11}&&&\\
 a_{21} &a_{22}&&\\
 \vdots &    \ddots  &\ddots&\\
  a_{n1}&    \cdots   & a_{n,n-1}    &a_{nn}\\
\end{vmatrix}=a_{11}a_{22}\cdots a_{nn}.\]
\end{example}
\pp

\begin{example}
\[\begin{vmatrix}
a_{11}&a_{12}&\cdots&a_{1n}\\
a_{21}& \iddots &\iddots&\\
\vdots&a_{n-1,2}& & \\
a_{n1}& & &\\
\end{vmatrix}=(-1)^{\frac{n(n-1)}{2}}a_{n1}a_{n-1,2}\cdots a_{1n}.\]
\end{example}
\end{frame}

\subsubsection{行列式行展开}
\begin{frame}\frametitle{行列式行展开}
\begin{theorem}[行列式行展开]
设$A=(a_{ij})_{n\times n}$为$n$阶方阵.则对任意$1\leq k \leq n$
\[\boxed{\det(A)=\sum_{j=1}^na_{kj}A_{kj} = \sum_{j=1}^{n}(-1)^{k+j}a_{kj}M_{kj}.}\]
\end{theorem}
\pp
\bigskip

\yang{思路: 不妨设$k=1$. 根据对行列式阶数的归纳, 将行列式
\[\det(A):=\sum_{i=1}^{n}(-1)^{i+1}a_{i1}M_{i1}\]
中除$M_{11}$以外的$M_{i1}$按第一行展开. \pp 同时右式
\[\sum_{i=1}^{n}(-1)^{i+j}a_{ij}M_{ij}\]
中除$M_{11}$以外的 $M_{1j}$ 按第一列展开. 然后比较两式是否相等即可.}
\end{frame}



\subsubsection{证明细节}
\begin{frame}\frametitle{证明细节}
\[ M\left[{{i_1i_2}\atop {j_1j_2}}\right]:=\det\left(A\left({1,\cdots,i_1-1,i_1+1,\cdots,i_2-1,i_2+1,\cdots,n} \atop {1,\cdots,j_1-1,j_1+1,\cdots,j_2-1,j_2+1,\cdots,n}\right)\right).\]
\pp
对行列式的阶数$n$进行归纳。当$n=1$,结论显然成立。假设结论对n-1阶行列式成立。\pp 则
\begin{equation*}
\begin{split}
\det(A)
& = \sum_{i=1}^{n}(-1)^{i+1}a_{i1}M_{i1} \\ \pp
& = a_{11}M_{11} + \sum_{i=2}^{n}(-1)^{i+1}a_{i1} \left(\sum_{j=2}^{n} (-1)^{1+(j-1)} a_{1j}\cdot M\left[{1i}\atop{1j}\right]\right)\\ \pp
& = a_{11}M_{11} + \sum_{j=2}^{n}(-1)^{1+j} a_{1j} \left(\sum_{i=2}^{n} (-1)^{(i-1)+1} a_{i1}\cdot M\left[{1i}\atop{1j}\right]\right)\\ \pp
& = \sum_{j=1}^{n} (-1)^{1+j}a_{1j} M_{1j}  \pp = \sum_{j=1}^{n} a_{1j}A_{1j} \\ \pp
\end{split}
\end{equation*}
\pp
其中第二个等式为对$n-1$阶行列式$M_{i1}$的第$1$行展开，第三个等式为交换求和号，第四个等式为对$n-1$阶行列式$M_{1j}$的第$1$列展开。
\end{frame}



\renewcommand{\mydate}{2025年03月20号\quad 第八次课}



\subsubsection{行列式基本性质}
\begin{frame}\frametitle{行列式基本性质}
\begin{theorem}
方阵$A$的行列式具有下列性质:
\begin{enumerate}
\item 交换$A$某两行(列)得$B$,则$\det(B)=-\det(A)$.\pp
\item 将$A$的某一行(列)乘$\lambda$得$B$,则$\det(B)=\lambda\det(A)$. \pp
\item 若$A$的某一行(列)是两个向量之和,则$\det(A)$可拆成两个行列式之和. \pp
\end{enumerate}
\end{theorem}
证明思路：将$\det(A)$按第$p，q$两行展开得
\[\det(A) = \sum_{i=1}^n \sum_{j=1}^{i-1} (-1)^{p+q+i+j-1}(a_{pi}a_{qj}-a_{pj}a_{qi}) M\left[{p,q}\atop{i,j}\right].\]

\pp

\yang{几何解释？}
\end{frame}









% ========== 2023年03月30号 第08次课 行列式 ==========
%\frame{\titlepage}






\begin{frame}\frametitle{行列式基本性质}
\begin{corollary}
\begin{enumerate}\setcounter{enumi}{3}
\item 若$A$的某两行(列)成比例,则$\det(A)=0$. \pp
\item 将$A$的某一行(列)$\lambda$倍加到另一行(列)得$B$,则$\det(B)=\det(A)$.
\end{enumerate}
\end{corollary}
\pp
\yang{几何解释？}
\pp

设 $A=(a_{ij})_{n\times n} = \begin{pmatrix}\alpha_1\\\alpha_2\\\vdots\\\alpha_n\end{pmatrix}$.则$\det$可以看成$\alpha_1,\cdots,\alpha_n$的函数. \pp  这一函数满足如下性质:
{\small
\begin{enumerate}
\item 反对称性: $\det(\cdots,\alpha_i,\cdots,\alpha_j\cdots)=-\det(\cdots,\alpha_j,\cdots,\alpha_i\cdots)$; \pp
\item 多重线性:  \\
$\det(\cdots,\lambda\alpha_i+\mu\beta_i,\cdots) = \lambda\det(\cdots,\alpha_i,\cdots) + \mu\det(\cdots,\beta_i,\cdots)$;
\pp
\item 规范性: $\det(\vec{e}_1,\vec{e}_2,\cdots,\vec{e}_n)=1$.
\end{enumerate}
}
\pp
\blue{注:} 这三条性质唯一地确定了函数$\det$. (几何意义！！)
\pp

\blue{例:} $\det(\alpha+\beta,\beta+\gamma,\gamma+\alpha)=2\det(\alpha,\beta,\gamma)$.
\end{frame}





\subsubsection{行列式显示表示}
\begin{frame}\frametitle{行列式显示表示}
\begin{equation*}
\begin{split}
\det(A)
& = \det\left(
\sum_{j_1=1}^{n}a_{1j_1}\vec{e}_{j_1}, \sum_{j_2=1}^{n}a_{2j_2}\vec{e}_{j_2},\cdots, \sum_{j_n=1}^{n}a_{nj_n}\vec{e}_{j_n}
\right)\\ \pp
& =\sum_{j_1=1}^{n}\sum_{j_2=1}^{n}\cdots \sum_{j_n=1}^{n} a_{1j_1}a_{2j_2}\cdots a_{nj_n}
\det\left(\vec{e}_{j_1}, \vec{e}_{j_2},\cdots, \vec{e}_{j_n}\right)\\ \pp
& =\sum_{(j_1j_2\cdots j_n)\in S_n} a_{1j_1}a_{2j_2}\cdots a_{nj_n}
\underbrace{\red{\boxed{\det\left(\vec{e}_{j_1}, \vec{e}_{j_2},\cdots, \vec{e}_{j_n}\right)}}}_{=??}\\
\end{split}
\end{equation*}
\pp
其中$S_n$为集合$\{1,\cdots,n\}$的全体排列集.
\end{frame}





\begin{frame}\frametitle{行列式显示表示}
\begin{definition}
\begin{enumerate}
\item 由$n$个两两不同的正整数组成的有序数组$(j_1,j_2,\cdots,j_n)$称为一个\blue{$n$元排列}. \pp  由$1,2,\cdots,n$组成的排列总数为$n!$. \pp
\item 称$(1,2,\cdots,n)$为\blue{标准排列}. \pp
\item 给定排列$(j_1,j_2,\cdots,j_n)$.若$1\leq p<q\leq n$且$j_p>j_q$,则称二元组$(j_p,j_q)$为排列$(j_1,j_2,\cdots,j_n)$的一个\blue{逆序}. \pp 例如: $(1,4,3,2)$的有三个逆序$(4,3),(4,2),(3,2)$. \pp
\item 排列$(j_1,j_2,\cdots,j_n)$的\blue{逆序总数}记为$\tau(j_1,j_2,\cdots,j_n)$. \pp  例如: $\tau(1,4,3,2)=3$. \pp
\item 若$\tau(j_1,j_2,\cdots,j_n)$为奇数,则称$(j_1,j_2,\cdots,j_n)$为\blue{奇排列}.否则称之为\blue{偶排列}. \pp
\item 将一个排列中的两个元素互换位置,其余元素位置不变.这个过程称为一个\blue{对换}.
\end{enumerate}
\end{definition}
\end{frame}






\begin{frame}\frametitle{排列与逆序}
\begin{lemma}
\begin{enumerate}
\item 对换改变奇偶性.\pp
\item 排列$(j_1,j_2,\cdots,j_n)$可经过$\tau(j_1,j_2,\cdots,j_n)$次相邻位置的对换变为标准排列.
\end{enumerate}
\end{lemma}
\pp
\begin{theorem}[行列式显示表示]
设$A=(a_{ij})_{n\times n}$为$n$阶方阵.则
\begin{equation*}
\det(A) = \sum_{(j_1j_2\cdots j_n)\in S_n}(-1)^{\tau(j_1,j_2,\cdots,j_n)} a_{1j_1}a_{2j_2}\cdots a_{nj_n}
\end{equation*}
\end{theorem}
\end{frame}





\subsubsection{二三阶行列式示意图}
\begin{frame}\frametitle{二三阶行列式示意图}
\begin{example}
\begin{enumerate}
\item $S_2=\{(12),(21)\}$, \pp $\tau(12)=0$, $\tau(21)=1$ \pp
\[\Rightarrow \begin{vmatrix}a&b\\c&d\end{vmatrix}=ad-bc.\]
\pp
\item $S_3=\{(123),(231),(312),(132),(213),(321)\}$,\pp  $\tau(123)=0$, $\tau(231)=\tau(312)=2$, $\tau(213)=\tau(132)=1$, $\tau(321)=3$,
\pp
\begin{itemize}
\item $\begin{vmatrix}
\red{a_{11}}&\blue{a_{12}}&\green{a_{13}}\\
\green{a_{21}}&\red{a_{22}}&\blue{a_{23}}\\
\blue{a_{31}}&\green{a_{32}}&\red{a_{33}}\\
\end{vmatrix}
{=\red{+a_{11}a_{22}a_{33}} \blue{+a_{12}a_{23}a_{31}} \green{+a_{13}a_{21}a_{32}} \atop \qquad -a_{11}a_{23}a_{32}-a_{12}a_{21}a_{33}-a_{13}a_{22}a_{31}}$
\pp
\item $\begin{vmatrix}
\red{a_{11}}&\blue{a_{12}}&\green{a_{13}}\\
\blue{a_{21}}&\green{a_{22}}&\red{a_{23}}\\
\green{a_{31}}&\red{a_{32}}&\blue{a_{33}}\\
\end{vmatrix}
{=+a_{11}a_{22}a_{33}+a_{12}a_{23}a_{31}+a_{13}a_{21}a_{32}\atop \qquad \red{-a_{11}a_{23}a_{32}} \blue{-a_{12}a_{21}a_{33}} \green{-a_{13}a_{22}a_{31}}}$
\end{itemize}
\end{enumerate}
\end{example}
\end{frame}





\begin{frame}
\begin{example} \footnotesize
\[\begin{vmatrix}
& & &a_{1}\\
& &\iddots&\\
&a_{n-1}& & \\
a_{n}& & &\\
\end{vmatrix}=(-1)^{\frac{(n)(n-1)}{2}}a_{1}a_{2}\cdots a_{n}.\]
\end{example}

\pp
\begin{example}
求下列三阶行列式: $\begin{vmatrix}
1 & 2 & 3\\
8 & 0 & 4\\
7 & 6 & 5
\end{vmatrix} \quad
\text{和}\quad \begin{vmatrix}
1 & 2 & 0\\
0 & 1 & 3\\
2 & 7 & 8
\end{vmatrix}$.
\end{example}
\end{frame}





\subsubsection{三角分块矩阵行列式}
\begin{frame}\frametitle{三角分块矩阵行列式}
\begin{proposition}
设$\small A=
\begin{pmatrix}
A_{11}&A_{12}&\cdots&A_{1k}\\
  &A_{22}&\cdots&A_{2k}\\
  &      &\ddots&\vdots\\
  &      &      &A_{kk}\\
\end{pmatrix}$为分块矩阵,其中$A_{ii}$均为方阵.则
\[\det(A)=\prod_{i=1}^k \det(A_{ii}).\]
\end{proposition}
\pp
{\small
证明思路：只需考虑$k=2$. 此时，不妨设拆分方式为 $n=r+(n-r)$.\pp 则}
\begin{equation*} \tiny
\begin{split}
\det(A) &= \sum_{(j_1,j_2,\cdots,j_n)\in S_n} (-1)^{\tau(j_1,\cdots,j_n)}a_{1j_1}a_{2j_2}\cdots a_{nj_n}\\ \pp
&=\sum_{(j_1,j_2,\cdots,j_r)\in S_r \leq r \atop(j_{r+1},\cdots,j_n)\in S_{n-r}(r+1,\cdots,n)} \left((-1)^{\tau(j_1,\cdots,j_r)}a_{1j_1}\cdots a_{r,j_r}\right) \cdot \left( (-1)^{\tau(j_{r+1},\cdots,j_n)}a_{r+1,j_{r+1}}\cdots a_{nj_n}\right)\\ \pp
%&=\left(\sum_{1\leq j_1,j_2,\cdots,j_r\leq r} (-1)^{\tau(j_1,\cdots,j_r)}a_{1j_1}\cdots a_{r,j_r}\right) \cdot \left(\sum_{r+1\leq j_{r+1},\cdots,j_n\leq n} (-1)^{\tau(j_{r+1},\cdots,j_n)}a_{r+1,j_{r+1}}\cdots a_{nj_n}\right) \\
&=\det(A_{11})\det(A_{22})\\
\end{split}
\end{equation*}
\end{frame}



\renewcommand{\mydate}{2025年03月25号\ 第九次课}

\subsubsection{行列式基本性质}
\begin{frame}\frametitle{行列式基本性质}
\begin{theorem}
\[\det(A)=\det(A^T).\]
\end{theorem}
\pp
证明：比较两边显示展开.(或，对阶数归纳，并利用展开式。)
\pp
\begin{theorem}
\[\det(AB)=\det(A)\det(B).\]
\end{theorem}
\pp
几何解释？
\pp
\begin{example}
设$m>n$, $A\in F^{m\times n}$, $B\in F^{n\times m}$. 则
\[\det(AB)=0.\]
\end{example}
\pp
\begin{example}
设$A\in F^{n\times n}$, $B\in F^{n\times m}$, $C\in F^{m\times n}$, $D\in F^{m\times m}$. 若$A$可逆,则
\[\det\begin{pmatrix} A&B\\C&D\\\end{pmatrix}=\det(A)\cdot\det(D-CA^{-1}B).\]
\end{example}
\end{frame}




\begin{frame}\frametitle{例}
\begin{example} 求
$\begin{vmatrix}
3& x + y + z & x^2 + y^2 + z^2 \\
x + y + z & x^2 + y^2 + z^2 & x^3 + y^3 + z^3 \\
x^2 + y^2 + z^2 & x^3 + y^3 + z^3 & x^4 + y^4 + z^4 \\
\end{vmatrix}$
\end{example}


\yang{
\[\begin{pmatrix}
3& x + y + z & x^2 + y^2 + z^2 \\
x + y + z & x^2 + y^2 + z^2 & x^3 + y^3 + z^3 \\
x^2 + y^2 + z^2 & x^3 + y^3 + z^3 & x^4 + y^4 + z^4 \\
\end{pmatrix} =
\begin{pmatrix}
1 & 1 & 1 \\
x & y & z \\
x^2 & y^2 & z^2 \\
\end{pmatrix}
\begin{pmatrix}
1 & x & x^2 \\
1 & y & y^2 \\
1 & z & z^2 \\
\end{pmatrix}\]}
\end{frame}

\subsection{可逆判定}


\subsubsection{伴随矩阵}
\begin{frame}\frametitle{伴随矩阵}
\pp
\begin{definition}[伴随矩阵]
设$A=(a_{ij})_{n\times n}$为$n$阶方阵.称矩阵
\[\small A^*:=\begin{pmatrix}
A_{11}&A_{21}&\cdots&A_{n1}\\
A_{12}&A_{22}&\cdots&A_{n2}\\
\cdots&\cdots&\cdots&\cdots\\
A_{1n}&A_{2n}&\cdots&A_{nn}\\
\end{pmatrix}=(A_{ij})_{n\times n}^T\]
为$A$的\blue{伴随矩阵}.
\end{definition}
\pp
\begin{theorem}
设$A=(a_{ij})_{n\times n}$为$n$阶方阵. 则
\[A^*A=AA^*=\det(A)I_n.\]
\end{theorem}
\pp
注：定理等价于 $\sum\limits_{k=1}^{n}a_{ik}A_{jk}=\det(A)\delta_{ij}$ 以及 $\sum\limits_{k=1}^{n}a_{ki}A_{kj}=\det(A)\delta_{ij}$, 其中$\delta_{ij}$为Kronecker记号,满足 $(\delta_{ij})_{n\times n}=I_n$.

\end{frame}





\subsubsection{矩阵可逆的判定条件}
\begin{frame}\frametitle{矩阵可逆的判定条件}
\begin{theorem}
设$A$为$n$阶方阵. 则以下几条等价 \pp
\begin{enumerate}
\item $A$可逆(i.e. 存在$X$使得$AX=I=XA$.); \pp
\item $\det(A)\neq0$; \pp
\item 存在$X$使得$I=XA$; \pp
\item 存在$X$使得$AX=I$.
\end{enumerate}
\pp
若以上几条成立,则$A^{-1}=\frac1{\det(A)}A^*$.
\end{theorem}
\end{frame}



\subsection{Cramer 法则}

\begin{frame}
\begin{theorem}[Cramer法则] \small
若$n$阶方阵$A=(a_{ij})_{n\times n}=(\alpha_1,\cdots,\alpha_n)$可逆.则一般线性方程组
\[AX=b\]
有唯一解
\[(x_1,\cdots,x_n)=\left(\frac{\Delta_1}{\Delta},\frac{\Delta_2}{\Delta},\cdots,\frac{\Delta_n}{\Delta}\right),\]
其中$\Delta=\det(A)$, $\Delta_i=\det(\alpha_1\cdots,\alpha_{i-1},b,\alpha_{i+1},\cdots,\alpha_n)$.
\end{theorem}
\pp
证明思路：$X=A^{-1}b=\frac{1}{\det(A)} A^*b$. \pp 几何解释？
\pp

理论层面上，这种方法既有用又美观，但在实际应用中，由于效率过低，通常不会被采用来解方程。
\pp
\begin{example}
使用Cramer法则求解线性方程组:
\[\footnotesize \left\{\begin{array}{rrrrr}
& + x_2 & + x_3 & + x_4 & = 1\\
x_1 &       & + x_3 & + x_4 & = 2\\
x_1 & + x_2 &       & + x_4 & = 3\\
x_1 & + x_2 & + x_3 &       & = 4\\
\end{array}\right. \pp \Rightarrow (x_1,x_2,x_3,x_4)=(7/3,4/3,1/3,-2/3).\]
\end{example}
\end{frame}


\renewcommand{\mydate}{2025年03月27号\ 第十次课}

\subsection{行列式的计算}


\subsubsection{行列式的计算}
\begin{frame}\frametitle{行列式的计算}
\begin{example}
$\begin{vmatrix}
0& 5&-4& 5\\
-3&-1&-5& 3\\
3& 1&-2&-3\\
-1& 4&-5&-1\\
\end{vmatrix}=665$.
\end{example}
\pp
\begin{example}
$\begin{vmatrix}
x&     1&\cdots& 1\\
1&     x&\ddots&\vdots\\
\vdots&\ddots&     x& 1\\
1&\cdots&     1& x\\
\end{vmatrix}=(x+n-1)(x-1)^{n-1}$.
\end{example}
\end{frame}





\begin{frame}
\begin{example} \small
$\begin{vmatrix}
x&      &      &      a_0\\
-1&\ddots&      &   \vdots\\
&\ddots&     x&  a_{n-2}\\
&      &    -1&x+a_{n-1}\\
\end{vmatrix}=x^n+a_{n-1}x^{n-1}+\cdots+a_1x+a_0$.
\end{example}
\pp
\begin{example} \small
$\begin{vmatrix}
1&     a_1&  \cdots&  a_1^{n-1}\\
1&     a_2&  \cdots&  a_2^{n-1}\\
\cdots&  \cdots&  \cdots&     \cdots\\
1&     a_n&  \cdots&  a_n^{n-1}\\
\end{vmatrix}=\prod\limits_{1\leq i<j\leq n} (a_j-a_i)$.
\end{example}
\pp
\begin{example} \small
$\Delta_n:=\begin{vmatrix}
2&     1&      &  \\
1&     2&\ddots&  \\
&\ddots&\ddots& 1\\
&      &     1& 2\\
\end{vmatrix}=n+1$.
\end{example}
\end{frame}


